% Clase 15 --- El movimiento circular y el período que no depende de v
% (§1.1--1.3).
% El panel (a) fija la notación que el docente estableció al empezar (punto =
% saliente, cruz = entrante) y muestra por qué el movimiento es circular. El (b)
% aisla el resultado que después hace posible el ciclotrón: al duplicar la
% velocidad se duplica el radio, así que el tiempo de vuelta es el mismo.
% Cuidado con el sentido: con B entrante ($-\hat z$) y q>0, la fuerza sobre una
% carga que va hacia la izquierda ($-\hat x$) apunta hacia abajo. Es la
% combinación que hace que el centro quede debajo de la partícula.
\begin{center}
\begin{tikzpicture}[scale=1]

  % =============== (a) por qué es circular ===============
  \begin{scope}
    \fill[figblue!5] (-2.3,-2.1) rectangle (2.3,2.1);
    \draw[figgray, line width=0.7pt] (-2.3,-2.1) rectangle (2.3,2.1);
    \foreach \x in {-1.9,-0.95,0,0.95,1.9} {
      \foreach \y in {-1.75,-0.9,-0.05,0.8,1.65} {
        \node[figblue!45, scale=0.6] at (\x,\y) {$\otimes$};
      }
    }
    \node[figlblaux, anchor=south west] at (-2.3,2.16)
      {$\vec B$ entrante ($\otimes$); saliente sería $\odot$};

    \draw[figaux, line width=0.9pt] (0,0) circle (1.2);
    \node[figcarga, fill=figred, minimum size=6pt] at (0,1.2) {};

    % velocidad hacia la izquierda y fuerza centrípeta (hacia abajo)
    \draw[figvecres] (-0.15,1.2) -- (-1.25,1.2);
    \node[figlbl, figred, anchor=south] at (-0.7,1.26) {$\vec v$};
    \draw[figvecaux, figblue, line width=1pt,
          -{Latex[length=2.2mm,width=1.6mm]}] (0,1.05) -- (0,0.25);
    \node[figlbl, figblue, anchor=west] at (0.1,0.65) {$\vec F$};

    \draw[figvecaux] (0,0) -- (300:1.2);
    \node[figlblaux, anchor=north east] at (300:0.62) {$r$};

    \node[figlbl, figblue, anchor=north, align=center] at (0,-2.45)
      {$\vec F = q\,\vec v\times\vec B$, siempre $\perp \vec v$};
    \node[figlbl, anchor=north, align=center] at (0,-3.2)
      {(a) módulo constante y $\perp\vec v$\\[-0.2em]
       \footnotesize $\Rightarrow$ circular uniforme, con
       $r = \dfrac{m\,v}{|q|\,B}$};
  \end{scope}

  % =============== (b) el período no depende de v ===============
  \begin{scope}[xshift=8.4cm]
    \fill[figblue!5] (-2.3,-2.1) rectangle (2.3,2.1);
    \draw[figgray, line width=0.7pt] (-2.3,-2.1) rectangle (2.3,2.1);
    \foreach \x in {-1.9,-0.95,0,0.95,1.9} {
      \foreach \y in {-1.75,-0.9,-0.05,0.8,1.65} {
        \node[figblue!45, scale=0.6] at (\x,\y) {$\otimes$};
      }
    }

    \draw[figaux, line width=0.9pt] (0,0.9) circle (0.75);
    \draw[figred, line width=0.9pt] (0,0.15) circle (1.5);
    \node[figcarga, fill=figred, minimum size=5pt] at (0,1.65) {};

    \node[figlblaux, anchor=east] at (-0.85,1.3) {$v$};
    \node[figlbl, figred, anchor=west] at (1.55,0.9) {$2v$};

    \node[figlbl, figred, anchor=north, align=center] at (0,-2.45)
      {mismo $T = \dfrac{2\pi m}{|q|\,B}$};
    \node[figlbl, anchor=north, align=center] at (0,-3.2)
      {(b) al duplicar $v$ se duplica $r$\\[-0.2em]
       \footnotesize $\Rightarrow$ el período no cambia};
  \end{scope}

\end{tikzpicture}\\[0.5em]
{\small\itshape El argumento del panel (a) se encadena en dos pasos que conviene
no saltear. Primero, como $\vec F \perp \vec v$ siempre, la fuerza no hace
trabajo y la \emph{rapidez} es constante; y como la fuerza está contenida en el
plano perpendicular a $\vec B$ y la velocidad inicial también, la partícula
nunca sale de ese plano. Recién entonces, con $v$ constante y
$\vec v \perp \vec B$ en todo instante, el módulo $|q|vB$ es constante: una
fuerza de módulo fijo y siempre perpendicular a la velocidad es exactamente la
receta del movimiento circular uniforme de Física I. El sentido de giro sale de
la regla de la mano derecha y se invierte si la carga es negativa, pero nada más
cambia. El panel (b) muestra la cancelación de $v$ que da
$T = 2\pi m/(|q|B)$: dos partículas con velocidades distintas recorren
circunferencias distintas y tardan lo mismo, porque la que va más rápido tiene
que recorrer proporcionalmente más camino. Ese hecho es lo que hace posible el
ciclotrón, y deja de valer cuando $v$ se acerca a la de la luz: ahí los efectos
relativistas hacen que $f$ dependa de la velocidad.}
\end{center}
